\chapter{$\kappa$-non-collapsed ancient solutions}\label{chap:kappa-ancient-solutions}
\label{kappasect}

\section{Preliminaries}

Throughout, $(M,g(t))$, $-\infty<t\leq0$, is a Ricci flow whose
time-slices are complete and have bounded nonnegative curvature operator.
In dimensions at most three this is equivalent to nonnegative sectional
curvature, and bounded curvature operator is equivalent to bounded scalar
curvature. Consequently $g(t)$ is non-increasing in $t$, while Hamilton's
Harnack inequality gives $\partial_tR\geq0$ (Corollary~\ref{posderiv};
see also \cite{Hamiltonharnack}).

\subsection{Definition}

\begin{definition}[$\kappa$-non-collapsed Ricci flow]\label{def:kappa-noncollapsed}
\uses{def:ricci-flow}
\label{kappadefn}

Fix $\kappa>0$ and $r_0>0$. A Ricci flow $(M,g(t))$, $a<t\leq b$,
of complete $n$-manifolds is \emph{$\kappa$-non-collapsed on scales at
most $r_0$} if, whenever $(p,t)\in M\times(a,b]$, $0<r\leq r_0$,
$a\leq t-r^2$, and
\[
 |\Rm(q,t')|\leq r^{-2}
 \quad(q\in B(p,t,r),\ t-r^2<t'\leq t),
\]
one has $\Vol B(p,t,r)\geq\kappa r^n$. It is
\emph{$\kappa$-non-collapsed} if this holds for every finite $r_0$.
\end{definition}

\begin{definition}[Ancient solution and $\kappa$-solution]\label{def:ancient-solution}
\uses{def:ricci-flow, def:kappa-noncollapsed}

An \emph{ancient solution} is a Ricci flow $(M,g(t))$,
$-\infty<t\leq0$, such that every time-slice is connected, complete,
non-flat, and has bounded nonnegative curvature operator. A
\emph{$\kappa$-solution} is an ancient solution that is
$\kappa$-non-collapsed on all scales.
\end{definition}

A $\kappa$-solution is also a $\kappa'$-solution for
$0<\kappa'\leq\kappa$.

\subsection{Examples}

\begin{example}[Shrinking round $2$-sphere is a $\kappa$-solution]\label{ex:shrinking-sphere-kappa-solution}
\uses{def:ancient-solution}

If $g_0$ is the round metric on $S^2$ with scalar curvature $1$, then
$g(t)=(1-t)g_0$ is an ancient Ricci flow. It is a $\kappa$-solution
whenever $\kappa$ does not exceed the volume of a unit ball in the
unit $2$-sphere.
\end{example}

\begin{theorem}[Two-dimensional $\kappa$-solutions are round]\label{thm:two-dimensional-kappa-solution-preview}
\uses{def:ancient-solution, cor:two-dimensional-gradient-shrinking-soliton}

Every orientable two-dimensional $\kappa$-solution is, after scaling,
the shrinking round $2$-sphere above.
\end{theorem}

\begin{example}[Round $n$-sphere and its quotients]\label{ex:round-sphere-quotient-kappa-solution}
\uses{def:ancient-solution}

Let $g_0$ be the round metric on $S^n$ with scalar curvature $n/2$.
Then $g(t)=(1-t)g_0$ is a $\kappa$-solution for every $\kappa$ no
larger than the volume of a unit ball in the unit $n$-sphere. If a
finite isometry group $\Gamma$ acts freely, the quotient flow on
$S^n/\Gamma$ is a $\kappa$-solution for every $\kappa$ no larger than
$|\Gamma|^{-1}$ times that volume.
\end{example}

\begin{example}[Shrinking cylinder $S^2\times\Ar$]\label{ex:cylinder-kappa-solution}
\uses{def:ancient-solution}

On $S^2\times\Ar$, the product metric
$g(t)=(1-t)g_0+ds^2$ is a $\kappa$-solution whenever $\kappa$ does
not exceed the volume of a unit ball in the product of the unit
$2$-sphere and $\Ar$.
\end{example}

\begin{example}[Quotient of the shrinking cylinder]\label{ex:cylinder-quotient-kappa-solution}
\uses{def:ancient-solution, ex:cylinder-kappa-solution}

The quotient of the shrinking cylinder by the product of the
antipodal map on $S^2$ and $s\mapsto-s$ on $\Ar$ is an orientable
$\kappa$-solution for some $\kappa>0$.
\end{example}

\begin{example}[Product with a circle is not $\kappa$-non-collapsed]\label{ex:cylinder-with-circle-not-kappa-noncollapsed}
\uses{def:ancient-solution, def:kappa-noncollapsed}

For the ancient product flow
$g(t)=(1-t)g_0+ds^2$ on $S^2\times S^1_R$,
\[
 |\Rm(p,t)|=(1-t)^{-1},\qquad
 \frac{\Vol_{g(t)}B(p,\sqrt{1-t})}{(1-t)^{3/2}}
 \leq\frac{8\pi^2R}{\sqrt{1-t}}\longrightarrow0
\]
as $t\to-\infty$. Hence it is not $\kappa$-non-collapsed for any
$\kappa>0$.
\end{example}

\subsection{A consequence of Hamilton's Harnack inequality}

\begin{proposition}[Harnack consequence: Lipschitz bound for reduced length]\label{prop:harnack-l-function-lipschitz-bound}
\uses{def:ricci-flow}
\label{uniformlip}

Let $(M,g(t))$, $-\tau_0\leq t\leq0$, be an $n$-dimensional Ricci
flow of complete manifolds with bounded nonnegative curvature
operator. Put $\tau=-t$, fix $x=(p,0)$, and let $0<c<1$. For
$\tau\leq(1-c)\tau_0$, on the full-measure open set where $l_x$ is
smooth,
\[
 |\nabla l_x|^2+R\leq\frac{(1+2c^{-1})l_x}{\tau},
 \qquad
 R-\frac{(1+c^{-1})l_x}{\tau}\leq\partial_\tau l_x.
\]
\end{proposition}
\begin{proof}
\sketch For
\[
 H(X)=-\partial_\tau R-\frac R\tau-2\langle\nabla R,X\rangle
       +2\Ric(X,X),
\]
Hamilton's Harnack inequality gives
$H(X)\geq-c^{-1}R/\tau$ on the stated interval. Along a minimizing
$\mathcal L$-geodesic $\gamma$ to $(q,\bar\tau)$ this yields
\begin{equation}\label{Kbartaueqn}
 K^{\bar\tau}(\gamma)
 =\int_0^{\bar\tau}\tau^{3/2}H(X)\,d\tau
 \geq-2c^{-1}\sqrt{\bar\tau}\,l_x(q,\bar\tau).
\end{equation}
Substitution in the two reduced-length identities of
Theorem~\ref{Ueqn} gives the asserted inequalities; the smooth-locus
inequalities extend as local $L^\infty$ inequalities.
\end{proof}

\begin{corollary}[Uniform reduced-length gradient and time-derivative bound]\label{cor:l-function-gradient-bound}
\uses{prop:harnack-l-function-lipschitz-bound}
\label{unianc}

For an ancient flow as above and $x=(p,0)$, every $\tau>0$ satisfies
\[
 |\nabla l_x|^2+R\leq\frac{3l_x}{\tau},
 \qquad
 -\frac{2l_x}{\tau}\leq\partial_\tau l_x
 \leq\frac{l_x}{\tau}.
\]
These hold smoothly where $l_x$ is smooth and as local
$L^\infty$ inequalities everywhere.
\end{corollary}
\begin{proof}
\sketch
\uses{prop:harnack-l-function-lipschitz-bound}
Let $\tau_0\to\infty$ and $c\to1$ in the proposition. This gives the
first inequality and the lower time-derivative bound. Concatenating
a minimizing $\mathcal L$-geodesic to $(q,\tau)$ with the stationary
path at $q$, differentiating its reduced length at $\tau$, and using
$R\leq3l_x/\tau$ gives the upper bound.
\end{proof}

\section{The asymptotic gradient shrinking soliton for
$\kappa$-solutions}\label{sect9.2}

Fix a $\kappa$-solution $(M,g(t))$, a point $x=(p,0)$, and points
$q(\tau)$ minimizing $l_x(\cdot,\tau)$, so that
$l_x(q(\tau),\tau)\leq n/2$. For $\bar\tau>0$ set
$g_{\bar\tau}(t)=\bar\tau^{-1}g(\bar\tau t)$.

\begin{theorem}[Existence of the asymptotic gradient shrinking soliton]\label{thm:asymptotic-gradient-shrinking-soliton}
\uses{def:ancient-solution, def:kappa-noncollapsed, def:gradient-shrinking-soliton}
\label{asympGSS}

Let $\bar\tau_k\to\infty$, put $g_k=g_{\bar\tau_k}$, and
$q_k=q(\bar\tau_k)$. After passage to a subsequence, the pointed
flows $(M,g_k(t),(q_k,-1))$, $-\infty<t<0$, converge smoothly to a
non-flat flow $(M_\infty,g_\infty(t),(q_\infty,-1))$. There is a
smooth $f:M_\infty\times(-\infty,0)\to\Ar$ such that
\begin{equation}\label{GSSequation}
 \Ric_{g_\infty(t)}+\Hess^{g_\infty(t)}f(t)
 +\frac1{2t}g_\infty(t)=0.
\end{equation}
The limit has nonnegative curvature operator, is
$\kappa$-non-collapsed, and satisfies $\partial_tR\geq0$.
\end{theorem}
\begin{proof}
\sketch The estimates and compactness argument are recorded in
Claims~\ref{claim:qk-length-bound}--\ref{cor:kappa-solution-geometric-limit-exists}.
The pulled-back reduced lengths converge to $l_\infty$; the
distributional argument culminating in
Proposition~\ref{prop:l-infinity-differential-equations} makes this
limit smooth. Proposition~\ref{prop:limit-gradient-shrinking-soliton-equation}
then gives Equation~\ref{GSSequation}. Finally,
Claim~\ref{claim:flat-limit-length-function-form} and the strict
reduced-volume bound rule out a flat limit.
\end{proof}

\begin{remark}[Caveats about the asymptotic soliton]\label{rem:asymptotic-soliton-caveats}
\uses{thm:asymptotic-gradient-shrinking-soliton}

The theorem does not assert bounded curvature on each limiting
time-slice, so in general the limit is not yet known to extend to a
$\kappa$-solution. Nor does it yet produce the usual self-similar
diffeomorphisms. Both conclusions are established below in dimensions
two and three (Corollaries~\ref{2DGSS} and~\ref{3DGSSkappa}).
\end{remark}

\subsection{Bounding the reduced length and curvature}

\begin{claim}[Reduced volume rescaling]\label{claim:reduced-volume-rescaling}
\uses{thm:asymptotic-gradient-shrinking-soliton}
\label{scaleinv2}

If $x_k=(p,0)$ in the $k$th rescaled flow, then
\[
 \widetilde V_{x_k}(\tau)=\widetilde V_x(\bar\tau_k\tau).
\]
\end{claim}
\begin{proof}
\sketch This is the reduced-volume reparameterization formula,
Lemma~\ref{Qformula}.
\end{proof}

\begin{corollary}[Reduced volume limit is constant]\label{cor:reduced-volume-limit-constant}
\uses{claim:reduced-volume-rescaling}
\label{redvolconst}

There is $V_\infty\in[0,(4\pi)^{n/2})$ such that, for every
$\tau>0$,
\begin{equation}\label{limVeqn}
 \lim_{k\to\infty}\widetilde V_{x_k}(\tau)=V_\infty.
\end{equation}
\end{corollary}
\begin{proof}
\sketch Reduced volume is positive and non-increasing, begins at
$(4\pi)^{n/2}$, and has a limit at infinity. Equality with
$(4\pi)^{n/2}$ would force the original flow to be flat by
Corollary~\ref{4picase}; combine this with
Claim~\ref{claim:reduced-volume-rescaling}.
\end{proof}

\begin{claim}[Bound for the reduced length at the minimizing point]\label{claim:qk-length-bound}
\uses{thm:asymptotic-gradient-shrinking-soliton, cor:l-function-gradient-bound}
\label{tausqtau}

For every $\tau>0$,
\[
 l_{x_k}(q_k,\tau)\leq\frac{n}{2\tau^2}+\frac{n\tau}{2}.
\]
\end{claim}
\begin{proof}
\sketch
\uses{cor:l-function-gradient-bound}
Scale invariance gives $l_{x_k}(q_k,1)\leq n/2$. Integrating the
lower bound $-2l/\tau\leq\partial_\tau l$ from $\tau$ to $1$ gives
$n/(2\tau^2)$ for $\tau<1$; integrating the upper bound
$\partial_\tau l\leq l/\tau$ gives $n\tau/2$ for $\tau>1$.
\end{proof}

\begin{corollary}[Quadratic distance bound for the reduced length]\label{cor:l-function-distance-bound}
\uses{claim:qk-length-bound}
\label{lxkest}

There is a positive continuous function $C_1(\tau)$ such that
\[
 l_{x_k}(q,\tau)\leq
 \left(\sqrt{\frac3\tau},d_{g_k(-\tau)}(q_k,q)+C_1(\tau)\right)^2,
\]
\[
 |\nabla l_{x_k}(q,\tau)|\leq
 \frac3\tau d_{g_k(-\tau)}(q_k,q)+
 \sqrt{\frac3\tau}C_1(\tau).
\]
\end{corollary}
\begin{proof}
\sketch
\uses{cor:l-function-gradient-bound,claim:qk-length-bound,prop:harnack-l-function-lipschitz-bound}
Integrate $|\nabla l|\leq\sqrt{3l/\tau}$ from $q_k$ to $q$, using
Claim~\ref{claim:qk-length-bound}; one may take
$C_1(\tau)=\sqrt{n/(2\tau^2)+n\tau/2}$. Substitute the resulting
bound into the gradient estimate.
\end{proof}

\begin{corollary}[Existence of the geometric limit flow]\label{cor:kappa-solution-geometric-limit-exists}
\uses{thm:asymptotic-gradient-shrinking-soliton, def:kappa-noncollapsed}

After passage to a subsequence, the flows
$(M,g_k(t),(q_k,-1))$ converge smoothly on $-\infty<t<0$ to a
complete flow $(M_\infty,g_\infty(t),(q_\infty,-1))$ with
nonnegative curvature operator. It is $\kappa$-non-collapsed on all
scales and satisfies $\partial_tR\geq0$.
\end{corollary}
\begin{proof}
\sketch The preceding bounds control $l_{x_k}$ and hence $R_{g_k}$
on every fixed based ball and every interval bounded away from
$t=0$. Nonnegative curvature controls $\Rm$, and non-collapsing plus
Theorem~\ref{flowlimit} gives a limit. Diagonalization in the final
time cutoff yields the full interval; curvature sign,
non-collapsing, and the Harnack consequence pass to the limit.
\end{proof}

\subsection{The limit function}

Using convergence embeddings on an exhaustion, pull $l_{x_k}$ and
$g_k$ back to functions $l_k$ and metrics $h_k$ on $M_\infty$.
The preceding estimates give local uniform Lipschitz bounds. After
passing to a subsequence, $l_k\to l_\infty$ strongly in
$C^{0,\alpha}_{\mathrm{loc}}$ for $0<\alpha<1$ and weakly in
$W^{1,2}_{\mathrm{loc}}$.

\begin{corollary}[Gradient bound for the limit length function]\label{cor:l-infinity-gradient-bound}
\uses{cor:l-function-distance-bound}
\label{linfest}

For $\tau>0$ and $q\in M_\infty$,
\[
 |\nabla l_\infty(q,\tau)|\leq
 \frac3\tau d_{g_\infty(-\tau)}(q_\infty,q)
 +\sqrt{\frac3\tau}C_1(\tau).
\]
\end{corollary}
\begin{proof}
\sketch Pass to the limit in
Corollary~\ref{cor:l-function-distance-bound}, using weak lower
semicontinuity.
\end{proof}

\begin{remark}[$l_\infty$ need not be a reduced length function]\label{rem:l-infinity-not-reduced-length}
The construction does not identify $l_\infty$ with reduced length
from any point of the limiting flow.
\end{remark}

\subsection{Differential inequalities for $l_\infty$}

\begin{proposition}[Differential equalities for the limit length function]\label{prop:l-infinity-differential-equations}
\uses{cor:l-infinity-gradient-bound}
\label{weakeqn}

The function $l_\infty$ is smooth on
$M_\infty\times(-\infty,0)$ and satisfies
\begin{equation}\label{weak1}
 \partial_\tau l_\infty+|\nabla l_\infty|^2-R
 +\frac n{2\tau}-\triangle l_\infty=0,
\end{equation}
\begin{equation}\label{weak2}
 2\triangle l_\infty-|\nabla l_\infty|^2+R
 +\frac{l_\infty-n}{\tau}=0.
\end{equation}
\end{proposition}
\begin{proof}
\sketch Lemma~\ref{lem:gradient-convergence-distributions} permits
passage to the weak reduced-length inequality. The weighted
distribution is nonnegative by
Corollary~\ref{cor:weighted-distribution-nonnegative}; the lower
bound in Lemma~\ref{lem:l-function-lower-bound} supplies the decay
needed to test against $1$. Constancy of reduced volume and
Claim~\ref{claim:laplacian-integral-vanishes} force the distribution
to vanish. Parabolic regularity gives smoothness and Equation
\ref{weak1}; the remaining reduced-length identity passes to the
limit and gives Equation~\ref{weak2}.
\end{proof}

\subsection{Preliminary results toward the proof of
Proposition~{\ref{prop:l-infinity-differential-equations}}}\label{sect9.2.3}

\begin{lemma}[Convergence of the gradient-squared term]\label{lem:gradient-convergence-distributions}
\uses{cor:l-infinity-gradient-bound}
\label{nablaconv}

For every $t<0$,
\[
 |\nabla^{h_k}l_k|_{h_k}^2d\vol(h_k)
 \longrightarrow |\nabla l_\infty|_{g_\infty}^2d\vol(g_\infty)
\]
in the distributional sense on $M_\infty\times\{t\}$.
\end{lemma}
\begin{proof}
\sketch Smooth convergence of $h_k$ and local gradient bounds give
\begin{equation}\label{firstdist}
 |\nabla^{h_k}l_k|_{h_k}^2d\vol(h_k)
 -|\nabla l_k|_{g_\infty}^2d\vol(g_\infty)\longrightarrow0.
\end{equation}
Weak lower semicontinuity yields one inequality. For a nonnegative
compactly supported $\varphi$, polarize the difference as
\begin{eqnarray}\label{11}
\int\varphi\bigl(|\nabla^{h_k}l_k|^2-|\nabla l_\infty|^2\bigr)d\vol(h_k)
&=&\int\langle\nabla^{h_k}(l_k-l_\infty),
\varphi\nabla^{h_k}l_k\rangle d\vol(h_k)\nonumber\\
&&+\int\langle\nabla^{h_k}(l_k-l_\infty),
\varphi\nabla l_\infty\rangle d\vol(h_k).
\end{eqnarray}
The second term tends to
\begin{equation}\label{12}
 \int\langle\nabla(l_k-l_\infty),
 \varphi\nabla l_\infty\rangle_{g_\infty}d\vol(g_\infty)=o(1)
\end{equation}
by weak $W^{1,2}$ convergence. The first has nonpositive limsup by
Claim~\ref{claim:nabla-convergence-limsup-bound}, giving the reverse
inequality.
\end{proof}

\begin{claim}[Vanishing limsup in the gradient identity]\label{claim:nabla-convergence-limsup-bound}
For positive $\epsilon_k\to0$ with
$l_\infty-l_k+\epsilon_k>0$ on $\supp\varphi$,
\[
 \limsup_{k\to\infty}\int
 \langle\nabla^{h_k}(l_k-l_\infty-\epsilon_k),
 \varphi\nabla^{h_k}l_k\rangle_{h_k}d\vol(h_k)\leq0.
\]
\end{claim}
\begin{proof}
\sketch The weak inequality from Theorem~\ref{weak}, first for
smooth nonnegative tests and then by $W^{1,2}$ approximation for
locally Lipschitz tests, is applied to
$l_\infty-l_k+\epsilon_k$. Its right-hand side tends to zero
uniformly on $\supp\varphi$; expanding the derivative of
$\varphi(l_k-l_\infty-\epsilon_k)$ leaves exactly the asserted term,
since the term involving $\nabla\varphi$ tends to zero.
\end{proof}

\begin{lemma}[The candidate distribution is non-positively paired]\label{lem:distribution-d-nonpositive}
\uses{lem:gradient-convergence-distributions}

Define
\[
 \mathcal D=\partial_\tau l_\infty+|\nabla l_\infty|^2-R
 +\frac n{2\tau}-\triangle l_\infty.
\]
It extends continuously to compactly supported $W^{1,2}$ functions.
For every compactly supported nonnegative Lipschitz function $\psi$,
$\mathcal D(\psi)\leq0$.
\end{lemma}
\begin{proof}
\sketch
\uses{lem:gradient-convergence-distributions}
Weak convergence of $\partial_\tau l_k$ and $\triangle^{h_k}l_k$,
together with Lemma~\ref{lem:gradient-convergence-distributions},
passes the weak reduced-length inequality to the limit:
\begin{equation}\label{1stlinftyeqn}
 \mathcal D\geq0
\end{equation}
in the source sign convention. Each zeroth-order term is locally
bounded, while the Laplacian acts by
$-\int\langle\nabla\psi,\nabla l_\infty\rangle$; hence the extension
to $W^{1,2}$. Approximation by nonnegative smooth functions gives the
stated pairing sign.
\end{proof}

\begin{corollary}[Weighted distribution is non-negative]\label{cor:weighted-distribution-nonnegative}
\uses{lem:distribution-d-nonpositive}
\label{etimesdis}

The map $\varphi\mapsto\mathcal D(e^{-l_\infty}\varphi)$ is a
distribution whose value on every nonnegative compactly supported
smooth $\varphi$ is nonnegative.
\end{corollary}
\begin{proof}
\sketch
\uses{lem:distribution-d-nonpositive}
The product $e^{-l_\infty}\varphi$ is nonnegative, compactly
supported, and Lipschitz, so the assertion is the preceding lemma
with the source sign convention.
\end{proof}

\subsection{Extension to non-compactly supported functions}

\begin{lemma}[Lower bound for the reduced length]\label{lem:l-function-lower-bound}
\uses{thm:asymptotic-gradient-shrinking-soliton}
\label{lowerbdd}

There is $c_1=c_1(n)>0$ such that, for $p,q\in M_k$,
\[
 l_{x_k}(p,\tau)\geq-l_{x_k}(q,\tau)-1
 +c_1\frac{d_{g_k(-\tau)}(p,q)^2}{\tau}.
\]
\end{lemma}
\begin{proof}
\sketch By scaling take $\tau=1$, choose minimizing
$\mathcal L$-geodesics $\gamma_1,\gamma_2$ to $p,q$, and set
$f(a,b,s)=d_{g_k(-s)}(a,b)$. Then
\begin{equation}\label{fform}
 d_{g_k(-1)}(p,q)=\int_0^1
 \bigl(\partial_sf+\langle\nabla_af,\gamma_1'\rangle
 +\langle\nabla_bf,\gamma_2'\rangle\bigr)\,ds.
\end{equation}
The gradient estimate and $\gamma_i'=\nabla l_{x_k}$ give
\begin{equation}\label{1eqn}
 |\langle\nabla_af,\gamma_1'\rangle|
 \leq\sqrt3s^{-3/4}\sqrt{l_{x_k}(p,1)+1},
\end{equation}
\begin{equation}\label{2eqn}
 |\langle\nabla_bf,\gamma_2'\rangle|
 \leq\sqrt3s^{-3/4}\sqrt{l_{x_k}(q,1)+1}.
\end{equation}
Local curvature bounds obtained from the same estimate and the
distance-variation inequality bound $\partial_sf$ by an integrable
multiple of these square roots. Integration gives
$d(p,q)\leq C(n)(\sqrt{l(p,1)+1}+\sqrt{l(q,1)+1})$; squaring and
rescaling proves the claim.
\end{proof}

\begin{corollary}[Refined lower bound for the reduced length]\label{cor:l-function-lower-bound-refined}
\uses{lem:l-function-lower-bound, claim:qk-length-bound}

For $q'\in M_k$ and $0<\tau_0\leq\tau'$,
\[
 l_{x_k}(q',\tau')\geq
 -\frac{n}{2(\tau')^2}-\frac{n\tau'}2-1
 +c_1\frac{d_{g_k(-\tau_0)}(q_k,q')^2}{\tau'}.
\]
\end{corollary}
\begin{proof}
\sketch
\uses{claim:qk-length-bound,lem:l-function-lower-bound}
Apply the lemma with $q=q_k$, use
Claim~\ref{claim:qk-length-bound}, and use monotonicity of the metric
to replace $d_{g_k(-\tau')}$ by $d_{g_k(-\tau_0)}$.
\end{proof}

\begin{claim}[Absolute convergence of the weighted distribution]\label{claim:absolute-convergence-weighted-distribution}
\uses{cor:l-function-lower-bound-refined}
\label{absconv1}

Fix a compact time interval $[-\tau,-\tau_0]$. If a locally
Lipschitz $f$ satisfies $|f|\leq C\max(l_\infty,1)$, then
$fe^{-l_\infty}$ is absolutely convergent: for every bounded smooth
$\varphi$ and every sequence $0\leq\psi_k\leq1$ of compactly
supported smooth functions that is eventually $1$ on each compact
set, the finite limit
\[
 \lim_{k\to\infty}(fe^{-l_\infty})(\varphi\psi_k)
\]
exists and is independent of the exhaustion.
\end{claim}
\begin{proof}
\sketch The refined lower bound gives
$l_\infty\geq c,d_{g_\infty(-\tau_0)}(q_\infty,\cdot)^2-C'$ outside
a fixed ball, uniformly on the time interval. Thus
$fe^{-l_\infty}$ decays exponentially, while nonnegative Ricci
curvature gives polynomial volume growth.
\end{proof}

\begin{corollary}[Absolute convergence of the gradient and time-derivative terms]\label{cor:absolute-convergence-gradient-terms}
\uses{claim:absolute-convergence-weighted-distribution, cor:l-function-gradient-bound}
\label{absconv}

The distributions $|\nabla l_\infty|^2e^{-l_\infty}$,
$Re^{-l_\infty}$, and
$|\partial_\tau l_\infty|e^{-l_\infty}$ are absolutely convergent
in the preceding sense.
\end{corollary}
\begin{proof}
\sketch
\uses{cor:l-function-gradient-bound,claim:absolute-convergence-weighted-distribution}
The reduced-length estimates bound all three coefficients by a
constant multiple of $\max(l_\infty,1)$ on each compact time
interval.
\end{proof}

\begin{claim}[Absolute convergence of the Laplacian term]\label{claim:absolute-convergence-laplacian-term}
\uses{claim:absolute-convergence-weighted-distribution}
\label{absconv2}

If $\varphi$ and the exhaustion $\psi_k$ above are uniformly
Lipschitz, then
\[
 \lim_{k\to\infty}\int_{M_\infty}
 \varphi\psi_k\triangle e^{-l_\infty}
 \,d\vol_{g_\infty}
\]
converges absolutely.
\end{claim}
\begin{proof}
\sketch Integrate by parts. The integrand becomes
$\langle\nabla(\varphi\psi_k),\nabla l_\infty\rangle e^{-l_\infty}$,
which is integrable because
$|\nabla l_\infty|\leq\max(1,|\nabla l_\infty|^2)$ and the previous
claim applies.
\end{proof}

\begin{corollary}[Non-negativity of the weak pairing]\label{cor:weak-inequality-nonnegative-integral}
\uses{claim:absolute-convergence-weighted-distribution, claim:absolute-convergence-laplacian-term, cor:weighted-distribution-nonnegative}
\label{weakneg}

If $0<\tau_0<\tau_1$, and $f$ is nonnegative, smooth, bounded, and
has bounded spatial gradient, then
\[
 \int_{\tau_0}^{\tau_1}\int_{M_\infty\times\{-\tau\}}
 \mathcal D f\tau^{-n/2}e^{-l_\infty}
 \,d\vol_{g_\infty}\,d\tau\geq0.
\]
\end{corollary}
\begin{proof}
\sketch
\uses{claim:absolute-convergence-weighted-distribution,claim:absolute-convergence-laplacian-term,cor:weighted-distribution-nonnegative}
Test Corollary~\ref{cor:weighted-distribution-nonnegative} against
$f\psi_k$, where $\psi_k$ are uniformly Lipschitz distance cutoffs.
The two absolute-convergence results justify passage to the limit.
\end{proof}

\subsection{Completion of the proof of
Proposition~{\ref{prop:l-infinity-differential-equations}}}\label{sect9.2.5}

\begin{claim}[The pure Laplacian integral vanishes]\label{claim:laplacian-integral-vanishes}
\uses{cor:absolute-convergence-gradient-terms}

For $0<\tau_0<\tau_1$,
\[
 \int_{\tau_0}^{\tau_1}\int_{M_\infty\times\{-\tau\}}
 \triangle e^{-l_\infty}\,d\vol_{g_\infty}\,d\tau
 =\int_{\tau_0}^{\tau_1}\int
 (|\nabla l_\infty|^2-\triangle l_\infty)e^{-l_\infty}
 \,d\vol_{g_\infty}\,d\tau=0.
\]
\end{claim}
\begin{proof}
\sketch
\uses{cor:absolute-convergence-gradient-terms}
Use uniformly Lipschitz cutoffs tending to $1$ and integrate by
parts; absolute convergence makes the boundary terms vanish.
\end{proof}

Constancy of $\widetilde V_\infty$ and the preceding claim give
\begin{equation}\label{tildeD}
 \int_{\tau_0}^{\tau_1}\int_{M_\infty\times\{-\tau\}}
 \mathcal D\tau^{-n/2}e^{-l_\infty}
 \,d\vol_{g_\infty}\,d\tau=0.
\end{equation}
Testing both $\varphi$ and $1-\varphi$ in the nonnegative weighted
distribution therefore gives
\begin{equation}\label{constcon}
 \mathcal D\tau^{-n/2}e^{-l_\infty}=0
\end{equation}
weakly. Parabolic regularity and the limiting reduced-length identity
give
\begin{equation}\label{wk=}
 2\triangle l_\infty-|\nabla l_\infty|^2+R
 +\frac{l_\infty-n}{\tau}=0.
\end{equation}

\subsection{The gradient shrinking soliton equation}

\begin{proposition}[The limit satisfies the gradient shrinking soliton equation]\label{prop:limit-gradient-shrinking-soliton-equation}
\uses{prop:l-infinity-differential-equations, def:gradient-shrinking-soliton}
\label{gradshrink}

On $M_\infty\times(-\infty,0)$, with $\tau=-t$,
\[
 \Ric_{g_\infty(t)}+\Hess^{g_\infty(t)}l_\infty(\cdot,\tau)
 -\frac1{2\tau}g_\infty(t)=0.
\]
\end{proposition}
\begin{proof}
  \uses{def:ricci-flow}
\sketch Apply Lemma~\ref{lem:conjugate-heat-equation-soliton-identity}
with $f=l_\infty$. Equation~\ref{weak1} says that $u$ satisfies the
conjugate heat equation, while Equation~\ref{weak2} makes the
corresponding $v$ identically zero. The evolution identity for $v$
then forces the displayed squared norm to vanish because $u>0$.
\end{proof}

\begin{lemma}[Conjugate heat equation identity for gradient shrinking solitons]\label{lem:conjugate-heat-equation-soliton-identity}
\uses{def:ricci-flow}
\label{9.1}

Let $(M,g(t))$, $0\leq t\leq T$, be an $n$-dimensional Ricci flow,
let $f$ be smooth, and put $\tau=T-t$. Then
\[
 u=(4\pi\tau)^{-n/2}e^{-f}
\]
satisfies $-\partial_tu-\triangle u+Ru=0$ if and only if
\[
 \partial_tf+\triangle f-|\nabla f|^2+R-\frac n{2\tau}=0.
\]
When this holds, the function
\[
 v=\bigl[\tau(2\triangle f-|\nabla f|^2+R)+f-n\bigr]u
\]
satisfies
\[
 -\partial_tv-\triangle v+Rv
 =-2\tau\left|\Ric_g+\Hess^g f-\frac1{2\tau}g\right|^2u.
\]
\end{lemma}
\begin{proof}
  \uses{claim:curvature-symmetries-bianchi}
\sketch Direct differentiation gives the equivalence. Set
$H=\tau(2\triangle f-|\nabla f|^2+R)+f-n$, so $v=Hu$. Then
\begin{equation}\label{dHeqn}
 \partial_tH=-(2\triangle f-|\nabla f|^2+R)+\partial_tf
 +\tau\bigl(2\partial_t\triangle f-
 \partial_t|\nabla f|^2+\partial_tR\bigr).
\end{equation}
Use Claims~\ref{claim:laplacian-time-derivative-commutation}
and~\ref{claim:laplacian-gradient-squared-formula}, together with
\begin{equation}\label{feqn}
 \partial_tf=-\triangle f+|\nabla f|^2-R+\frac n{2\tau},
\end{equation}
and $\partial_tR=\triangle R+2|\Ric|^2$. Completing the square gives
\[
 \partial_tH+\triangle H+
 \frac{2\langle\nabla u,\nabla H\rangle}{u}
 =2\tau\left|\Ric+\Hess f-\frac1{2\tau}g\right|^2.
\]
The product rule for $v=Hu$ yields the result.
\end{proof}

\begin{claim}[Commuting the Laplacian and the time derivative of $f$]\label{claim:laplacian-time-derivative-commutation}
\label{11.20}
Along Ricci flow,
\[
 \partial_t\triangle f=\triangle(\partial_tf)
 +2\langle\Ric,\Hess f\rangle.
\]
\end{claim}
\begin{proof}
\sketch
\uses{claim:curvature-symmetries-bianchi}
Differentiate $g^{ij}\Hess(f)_{ij}$. Besides the two displayed
terms, this produces
$-g^{ij}(\partial_t\Gamma^k_{ij})\partial_kf$. In normal
coordinates,
\[
 g^{ij}\partial_t\Gamma^k_{ij}
 =g^{kl}g^{ij}(-\nabla_j\Ric_{li}-\nabla_i\Ric_{lj}
 +\nabla_l\Ric_{ij})=0
\]
by the contracted second Bianchi identity.
\end{proof}

\begin{claim}[Bochner formula for $|\nabla f|^2$]\label{claim:laplacian-gradient-squared-formula}
\uses{lem:laplacian-one-form}

\[
 \triangle|\nabla f|^2
 =2\langle\nabla\triangle f,\nabla f\rangle
 +2\Ric(\nabla f,\nabla f)+2|\Hess f|^2.
\]
\end{claim}
\begin{proof}
\sketch
\uses{lem:laplacian-one-form}
Apply the product rule to $\langle df,df\rangle$ and use
$\triangle df=d\triangle f+\Ric(\nabla f,\cdot)$.
\end{proof}

\subsection{Completion of the proof of
Theorem~{\ref{thm:asymptotic-gradient-shrinking-soliton}}}

\begin{claim}[The pullback of $l_\infty$ in the flat case]\label{claim:flat-limit-length-function-form}
\uses{prop:limit-gradient-shrinking-soliton-equation}

If $(M_\infty,g_\infty(t))$ is flat for some $t<0$, then it is
isometric to $\Ar^n$, and under such an isometry
\[
 l_\infty(x,\tau)=\frac{|x|^2}{4\tau}+\langle x,a\rangle+b\tau
\]
for some $a\in\Ar^n$ and $b\in\Ar$.
\end{claim}
\begin{proof}
\sketch
\uses{lem:conjugate-heat-equation-soliton-identity,prop:limit-gradient-shrinking-soliton-equation}
On the Euclidean universal cover the soliton equation reads
$\Hess l_\infty=g/(2\tau)$, so the difference from
$|x|^2/(4\tau)$ is affine. Such a function is invariant under no
nontrivial free deck action, hence the covering is trivial. The
conjugate heat equation fixes the time-dependent constant in the
stated form.
\end{proof}

If the limit were flat, the claim and Theorem~\ref{4picase} would
give $\widetilde V_\infty=(4\pi)^{n/2}$, contradicting
Corollary~\ref{cor:reduced-volume-limit-constant}. This completes the
non-flatness part of Theorem~\ref{thm:asymptotic-gradient-shrinking-soliton}.

\section{Splitting results at infinity}

\subsection{Point-picking}

\begin{lemma}[Point-picking lemma]\label{lem:point-picking-general}
\label{4times}
Let $(M,g)$ be Riemannian, let $B(p,2r)$ have compact closure, and
let $f:B(p,2r)\times(-2r,0]\to\Ar$ be continuous and bounded with
$f(p,0)>0$. There is $(q,t)$ in this cylinder such that, for
$\alpha=f(p,0)/f(q,t)$,
\begin{enumerate}
\item[(1)] $f(q,t)\geq f(p,0)$;
\item[(2)] $d(p,q)\leq2r(1-\alpha)$ and
$t\geq-2r(1-\alpha)$;
\item[(3)] $f(q',t')<2f(q,t)$ on
$B(q,\alpha r)\times(t-\alpha r,t]$.
\end{enumerate}
\end{lemma}
\begin{proof}
\sketch Starting at $(p,0)$, whenever (3) fails choose a point in
the indicated parabolic neighborhood where $f$ at least doubles.
After $j$ steps the accumulated spatial and backward-time
displacements are at most
$r(1+2^{-1}+\cdots+2^{1-j})<2r$, whereas $f\geq2^jf(p,0)$.
Boundedness of $f$ forces termination, and the geometric sums give
(1)--(3).
\end{proof}

\begin{corollary}[Time-independent point-picking]\label{cor:point-picking-time-independent}
\uses{lem:point-picking-general}
\label{pp}

If $f:B(p,2r)\to\Ar$ is continuous, bounded, and positive at $p$,
then some $q\in B(p,2r)$ satisfies $f(q)\geq f(p)$ and, with
$\alpha=f(p)/f(q)$,
\[
 d(p,q)\leq2r(1-\alpha),\qquad
 f(q')<2f(q)\quad(q'\in B(q,\alpha r)).
\]
\end{corollary}
\begin{proof}
\sketch
\uses{lem:point-picking-general}
Apply the lemma to the time-independent extension of $f$.
\end{proof}

\subsection{Splitting results}

\begin{proposition}[Splitting at infinity]\label{prop:splitting-at-infinity}
\uses{def:kappa-noncollapsed}
\label{prodatinf}

Let $(M,g(t))$, $-\infty<t<0$, be a $\kappa$-non-collapsed flow of
dimension $n\leq3$ whose time-slices are complete, noncompact,
non-flat, and have nonnegative curvature operator. Assume
$\partial_tR\geq0$. If, for some $p\in M$, points $p_i\to\infty$
satisfy
\[
 R(p_i,-1)d_{g(-1)}(p,p_i)^2\longrightarrow\infty,
\]
then there are $q_i\to\infty$, with $Q_i=R(q_i,-1)$ and
$Q_id_{g(-1)}(p,q_i)^2\to\infty$, such that
\[
 (M,Q_ig(Q_i^{-1}(t+1)-1),(q_i,-1))
 \longrightarrow (N^{n-1},h(t))\times(\Ar,ds^2).
\]
The product limit exists for $-\infty<t\leq-1$, and
$(N,h(-1))$ is non-flat with bounded nonnegative curvature.
The splitting assertion is valid in every dimension, although only
the cases $n\leq3$ are proved and used here.
\end{proposition}
\begin{proof}
\sketch
\uses{cor:point-picking-time-independent}
Apply Corollary~\ref{cor:point-picking-time-independent} to
$\sqrt R$ on $B(p_i,d_i/2)$, $d_i=d(p,p_i)$. This gives $q_i$ and
balls whose rescaled radii tend to infinity and on which
$R(\cdot,-1)\leq4Q_i$. The hypothesis $\partial_tR\geq0$ extends
this bound backward in time. Non-collapsing and
Theorem~\ref{flowlimit} give a complete limit with
$R(q_\infty,-1)=1$. Theorem~\ref{topsplit} splits its $-1$ slice
off a line; Corollary~\ref{localprod} propagates the splitting
through the flow. The remaining factor is non-flat and has bounded
nonnegative curvature.
\end{proof}

\begin{corollary}[No two-dimensional example]\label{cor:no-two-dimensional-splitting-example}
\uses{prop:splitting-at-infinity}
\label{compact2D}

No two-dimensional flow satisfies the hypotheses of
Proposition~\ref{prop:splitting-at-infinity}.
\end{corollary}
\begin{proof}
  \uses{cor:point-picking-time-independent}
\sketch The proposition would produce a product of two
one-dimensional manifolds, hence a flat limit, contradicting the
normalization $R(q_\infty,-1)=1$.
\end{proof}

\section{Classification of gradient shrinking solitons in dimensions $2$ and $3$}

At time $-1$ the soliton equation is
\begin{equation}\label{GSS-1}
 \Ric_{g_\infty(-1)}+\Hess^{g_\infty(-1)}f
 -\frac12g_\infty(-1)=0.
\end{equation}

The noncompact model is
$g(t)=|t|h_{-1}+ds^2$ on $S^2\times\Ar$, where $h_{-1}$ has
Gaussian curvature $1/2$, with potential $f=s^2/(4|t|)$. Its
quotient by the product of $s\mapsto-s$ and the antipodal map on
$S^2$ is again a gradient shrinking soliton.

\begin{definition}[$\kappa$-non-collapsed Riemannian manifold]\label{def:kappa-noncollapsed-riemannian-manifold}
\uses{def:kappa-noncollapsed}

A complete $n$-manifold $(M,g)$ is \emph{$\kappa$-non-collapsed}
if, for every $p\in M$ and $r>0$, the bound
$|\Rm_g|\leq r^{-2}$ on $B(p,r)$ implies
$\Vol B(p,r)\geq\kappa r^n$.
\end{definition}

\begin{theorem}[Classification of gradient shrinking solitons in dimension 2 and 3]\label{thm:gradient-shrinking-soliton-classification}
\uses{def:kappa-noncollapsed-riemannian-manifold, def:gradient-shrinking-soliton}
\label{GSSclass}

Let $(M,g)$ be a complete non-flat $\kappa$-non-collapsed manifold
of dimension two or three with bounded nonnegative curvature. If a
$C^2$ function $f$ satisfies
\[
 \Ric_g+\Hess^g f=\frac12g,
\]
then $g$ extends to a Ricci flow $(M,G(t))$, $-\infty<t<0$, with
$G(-1)=g$ and $(M,G(t))$ isometric to $(M,|t|g)$. The flow is exactly
one of:
\begin{enumerate}
\item[(1)] a shrinking family of compact round manifolds;
\item[(2)] a product of shrinking round $2$-spheres with $\Ar$;
\item[(3)] the quotient of the flow in (2) by an isometric involution.
\end{enumerate}
\end{theorem}
\begin{proof}
\sketch Boundedness of $\Hess f$ makes the gradient flow $\Phi_s$
complete. The self-similar extension
\begin{equation}\label{Gdefn}
 G(t)=|t|\Phi_{-\log|t|}^*g,\qquad -\infty<t<0,
\end{equation}
satisfies Ricci flow. Claim~\ref{claim:compact-positive-curvature-round}
handles the compact positive-curvature case;
Claim~\ref{claim:non-strictly-positive-curvature-splitting} reduces
the non-strict case to a surface times a line; and
Proposition~\ref{prop:no-noncompact-positive-curvature-soliton}
excludes the remaining case. The two-dimensional classification
then makes the surface factor round, yielding (1)--(3).
\end{proof}

\subsection{Compact positive curvature}\label{compactcs}

\begin{claim}[Compact positive-curvature case is round]\label{claim:compact-positive-curvature-round}
\uses{thm:gradient-shrinking-soliton-classification}

If $(M,g,f)$ satisfies the classification theorem and $M$ is
compact with positive curvature, then the flow in
Equation~\ref{Gdefn} is a shrinking family of round manifolds.
\end{claim}
\begin{proof}
\sketch Self-similarity makes every scale- and
diffeomorphism-invariant curvature ratio constant in time. In
dimension three Hamilton's pinching theorem
(Theorem~\ref{flowtoround}) forces that ratio to tend to one, hence
it is already one. In dimension two Hamilton's surface result
implies the same conclusion \cite{Hamiltonsurface}; see also
\cite{ChowKnopf}.
\end{proof}

\begin{corollary}[Compact asymptotic soliton forces time-shifted equivalence]\label{cor:compact-asymptotic-soliton-classification}
\uses{claim:compact-positive-curvature-round}
\label{compactcase}

If a three-dimensional $\kappa$-solution has a compact asymptotic
gradient shrinking soliton, then it is isomorphic to a time shift of
that soliton.
\end{corollary}
\begin{proof}
\sketch
\uses{claim:compact-positive-curvature-round}
The asymptotic soliton is round. Smooth convergence makes the
rescaled original slices increasingly pinched; scale invariance and
Hamilton--Ivey pinching force every original slice to be round. A
round flow is determined, up to isomorphism, by its manifold and
singular time, so only a time shift remains.
\end{proof}

\subsection{Non-strictly positive curvature}

\begin{claim}[Non-strictly positive curvature splits off a line]\label{claim:non-strictly-positive-curvature-splitting}
\uses{thm:gradient-shrinking-soliton-classification}
\label{flatssplit}

If $(M,g,f)$ satisfies the classification theorem but $(M,g)$ is not
strictly positively curved, then $n=3$ and the self-similar flow has
a one- or two-sheeted cover equal to the product of a positively
curved, $\kappa$-non-collapsed surface flow of bounded curvature and
a fixed flat line.
\end{claim}
\begin{proof}
\sketch Hamilton's strong maximum principle gives a one- or
two-sheeted product cover with a flat one-dimensional factor. For a
unit vector $Y$ in that factor, Equation~\ref{GSS-1} gives
$\Hess\widetilde f(Y,Y)=1/2$ (equivalently the source normalization
states this as a fixed positive constant), excluding a circle
factor. Non-collapsing and bounded curvature pass to the surface
factor.
\end{proof}

\subsection{Noncompact strict positive curvature}\label{noncomp}

\begin{proposition}[No non-compact positively curved examples]\label{prop:no-noncompact-positive-curvature-soliton}
\uses{thm:gradient-shrinking-soliton-classification}
\label{noncompsol}

No two- or three-dimensional soliton satisfying the classification
theorem is both noncompact and positively curved.
\end{proposition}
\begin{proof}
\sketch Tracing and taking divergence of the soliton equation gives
\begin{equation}\label{Rfeqn}
 dR+\triangle(df)-\Ric(\nabla f,\cdot)=0,
\end{equation}
and the contracted Bianchi identity yields
\begin{equation}\label{dR}
 dR=2\Ric(\nabla f,\cdot).
\end{equation}
Claims~\ref{claim:ricci-integral-bound}
and~\ref{claim:ricci-cauchy-schwarz-bound} imply, along every
minimizing ray $\gamma$, that
\begin{equation}\label{sqrts}
 \int_0^{\bar s}|\Ric(\gamma',\widetilde E_i)|\,ds
 \leq C'\sqrt{\bar s},
\end{equation}
\begin{equation}\label{X}
 \gamma'(f)(\bar s)\geq\frac{\bar s}{2}-C'',
\end{equation}
\begin{equation}\label{sqrts2}
 |\widetilde E_i(f)(\gamma(\bar s))|
 \leq C''(\sqrt{\bar s}+1).
\end{equation}
Thus $f$ is proper, its gradient is asymptotically radial, and
Equation~\ref{dR} makes $R$ nondecreasing on gradient lines.
Splitting at infinity excludes dimension two and gives in dimension
three a round shrinking cylinder limit. Its normalization is fixed
by Claim~\ref{claim:cylinder-limit-scalar-curvature-normalization}.
Large level spheres of $f$ then have area below $8\pi$ and, by the
Gauss equation,
\begin{equation}\label{R_N}
 R_N\leq R-2\Ric(e_3,e_3)
 +\frac{(1-R+\Ric(e_3,e_3))^2}{2|\nabla f|^2}<1.
\end{equation}
This contradicts Gauss--Bonnet.
\end{proof}

\begin{claim}[Bounded Ricci integral along minimizing rays]\label{claim:ricci-integral-bound}
There is a constant $C$, independent of a unit-speed minimizing
geodesic $\gamma:[0,\bar s]\to M$ from a fixed point and of
$\bar s$, such that
\[
 \int_0^{\bar s}\Ric(\gamma',\gamma')\,ds\leq C.
\]
\end{claim}
\begin{proof}
\sketch Use the second variation with parallel orthonormal fields
multiplied by cutoff functions that rise linearly on $[0,1]$, equal
one on $[1,\bar s-1]$, and fall linearly at the other end:
\begin{equation}\label{Yeqn}
 \int_0^{\bar s}-\langle\mathcal R(Y,\gamma')Y,\gamma'\rangle ds
 \leq\int_0^{\bar s}|\nabla_{\gamma'}Y|^2ds.
\end{equation}
Summing over the transverse fields bounds the middle Ricci integral
by $2(n-1)$; bounded curvature controls both endpoint pieces.
\end{proof}

\begin{claim}[Cauchy--Schwarz bound for the Ricci tensor]\label{claim:ricci-cauchy-schwarz-bound}
For every unit vector $X$,
\[
 |\Ric(X,\cdot)|^2\leq R\Ric(X,X).
\]
\end{claim}
\begin{proof}
\sketch Diagonalize the nonnegative Ricci tensor with eigenvalues
$\lambda_i$. Then the two sides are
$\sum_i(X^i)^2\lambda_i^2$ and
$(\sum_i\lambda_i)\sum_i(X^i)^2\lambda_i$, respectively.
\end{proof}

\begin{corollary}[Summary: non-compact soliton splitting]\label{cor:noncompact-soliton-splitting-summary}
\uses{claim:non-strictly-positive-curvature-splitting, prop:splitting-at-infinity}
\label{compact2Dsol}

There is no noncompact two-dimensional manifold satisfying the
classification theorem. For a noncompact positively curved
three-dimensional example, points tending to infinity and
approaching the supremum of $R$ admit a based flow limit that is a
product $(\Sigma^2,h(t))\times\Ar$, where the surface flow is
$\kappa$-non-collapsed and has positive bounded curvature.
\end{corollary}
\begin{proof}
\sketch Non-strict curvature is handled by
Claim~\ref{claim:non-strictly-positive-curvature-splitting}; in the
strict case the estimates in the preceding proposition meet the
hypotheses of Proposition~\ref{prop:splitting-at-infinity}. Its
two-dimensional conclusion is impossible, and its
three-dimensional limit has the asserted form.
\end{proof}

\begin{corollary}[Classification of two-dimensional gradient shrinking solitons]\label{cor:two-dimensional-gradient-shrinking-soliton}
\uses{cor:no-two-dimensional-splitting-example, cor:noncompact-soliton-splitting-summary}
\label{2DGSS}

\begin{enumerate}
\item[(1)] A two-dimensional flow satisfying all hypotheses of
Proposition~\ref{prop:splitting-at-infinity} except possibly
noncompactness is compact; after restricting to $(-\infty,-a]$ and
shifting time by $a>0$, it is a $\kappa$-solution.
\item[(2)] Every asymptotic gradient shrinking soliton of a
two-dimensional $\kappa$-solution is a shrinking family of compact
round surfaces.
\item[(3)] Every two-dimensional $\kappa$-solution is a shrinking
family of compact round surfaces.
\end{enumerate}
\end{corollary}
\begin{proof}
\sketch
\uses{cor:no-two-dimensional-splitting-example,cor:noncompact-soliton-splitting-summary,thm:asymptotic-gradient-shrinking-soliton}
The first statement follows from
Corollary~\ref{cor:no-two-dimensional-splitting-example}. An
unbounded-curvature asymptotic soliton would satisfy the splitting
hypotheses, so it has bounded curvature and is compact by
Corollary~\ref{cor:noncompact-soliton-splitting-summary}; Hamilton's
surface classification makes it round \cite{Hamiltonsurface}; see
also \cite{ChowKnopf}. For a $\kappa$-solution, both the backward
asymptotic limit and the forward singular limit are round
\cite{ChowKnopf}. Monotonicity of entropy, with
strictness away from shrinking solitons, forces the entire flow to
be round \cite{ChowKnopf}.
\end{proof}

\begin{claim}[Normalization of the limiting cylinder's curvature]\label{claim:cylinder-limit-scalar-curvature-normalization}
For the cylinder obtained at infinity in the positively curved
three-dimensional case, the scalar curvature of $(S^2,h(-1))$ is
$1$.
\end{claim}
\begin{proof}
\sketch Existence for all $t<0$ gives the upper bound $1$. Since
$R$ is increasing along gradient lines, $\inf R_g>0$ and the
self-similar flow has $\inf R_{G(t)}=(\inf R_g)/|t|\to\infty$ as
$t\uparrow0$. Hence the limiting sphere becomes singular at zero,
forcing equality.
\end{proof}

\subsection{Completion of the classification}

\begin{corollary}[The splitting-at-infinity limit is a family of round surfaces]\label{cor:three-dimensional-splitting-compact-surfaces}
\uses{prop:splitting-at-infinity, cor:two-dimensional-gradient-shrinking-soliton}
\label{3Dinfty}

For a three-dimensional flow satisfying
Proposition~\ref{prop:splitting-at-infinity}, its splitting limit is
the product of a shrinking family of compact round surfaces with a
line. This applies in particular to noncompact asymptotic solitons
of three-dimensional $\kappa$-solutions.
\end{corollary}
\begin{proof}
\sketch
\uses{prop:splitting-at-infinity,cor:two-dimensional-gradient-shrinking-soliton}
The surface factor inherits nonflat nonnegative bounded curvature,
$\partial_tR\geq0$, and $\kappa'$-non-collapsing. The
two-dimensional classification therefore makes it compact and
round.
\end{proof}

\begin{corollary}[Asymptotic soliton of a 3D $\kappa$-solution has bounded curvature]\label{cor:three-dimensional-asymptotic-soliton-bounded-curvature}
\uses{cor:three-dimensional-splitting-compact-surfaces, cor:two-dimensional-gradient-shrinking-soliton}
\label{3DGSSkappa}

Every time-slice of an asymptotic gradient shrinking soliton of a
three-dimensional $\kappa$-solution has bounded curvature.
Consequently, restriction to $t\leq-a$ followed by a time shift by
$a>0$ is a $\kappa$-solution.
\end{corollary}
\begin{proof}
\sketch
\uses{cor:three-dimensional-splitting-compact-surfaces,cor:two-dimensional-gradient-shrinking-soliton}
If curvature is not strictly positive, the maximum principle gives
a finite cover splitting as a round surface flow times a line. If
positive curvature were unbounded, point-picking and
Corollary~\ref{cor:three-dimensional-splitting-compact-surfaces}
would produce arbitrarily small cylindrical necks. The soul theorem
makes the manifold $\Ar^3$, so the surface is $S^2$; this
contradicts Proposition~\ref{narrows}.
\end{proof}

\begin{corollary}[The asymptotic soliton falls under the classification theorem]\label{cor:asymptotic-soliton-classification-application}
\uses{cor:three-dimensional-asymptotic-soliton-bounded-curvature, thm:gradient-shrinking-soliton-classification, thm:asymptotic-gradient-shrinking-soliton}

Every asymptotic gradient shrinking soliton of a three-dimensional
$\kappa$-solution is one of the three flows in
Theorem~\ref{thm:gradient-shrinking-soliton-classification}.
\end{corollary}
\begin{proof}
\sketch
\uses{cor:three-dimensional-asymptotic-soliton-bounded-curvature,thm:gradient-shrinking-soliton-classification,thm:asymptotic-gradient-shrinking-soliton}
Bounded curvature and $\partial_tR\geq0$ convert parabolic
non-collapsing into $\kappa'$-non-collapsing of the $-1$ slice.
The soliton equation therefore meets the classification theorem.
Uniqueness identifies the classified self-similar flow with the
original limit: directly in the compact case, on the surface factor
in the product case, and after passage to the finite cover in the
quotient case.
\end{proof}

\subsection{Asymptotic curvature}

\begin{definition}[Asymptotic scalar curvature]\label{def:asymptotic-scalar-curvature}
For a complete connected noncompact manifold $(M,g)$ of
nonnegative curvature and $p\in M$, define
\[
 \mathcal R(M,g)=\limsup_{x\to\infty}R(x)d_g(x,p)^2.
\]
This is independent of $p$.
\end{definition}

\begin{proposition}[Non-compact $\kappa$-solutions have infinite asymptotic curvature]\label{prop:asymptotic-scalar-curvature-infinite}
\uses{def:asymptotic-scalar-curvature, def:ancient-solution, def:kappa-noncollapsed}
\label{asympscal}

If $(M,g(t))$, $-\infty<t<0$, is a connected noncompact
$\kappa$-solution of dimension at most three, then
$\mathcal R(M,g(t))=+\infty$ for every $t<0$.
The conclusion is valid in every dimension; the present proof uses
the low-dimensional results established in this chapter.
\end{proposition}
\begin{proof}
\sketch
\uses{cor:two-dimensional-gradient-shrinking-soliton,claim:non-strictly-positive-curvature-splitting}
The two-dimensional assertion is vacuous. In dimension three, a
non-strictly positively curved solution has a finite cover splitting
off a line, so its asymptotic curvature is infinite. Assume strict
positive curvature. If $0<\mathcal R<\infty$, rescale about points
$x_n\to\infty$ by $Q_n=R(x_n,t)$ with
$Q_nd(x_n,p)^2\to\mathcal R$. Curvature control on annuli, Shi
estimates, and non-collapsing produce a nonflat local ancient smooth
limit. The same rescalings converge in the Gromov--Hausdorff sense
to a rescaled Tits cone, contradicting Proposition~\ref{nocones}.

If $\mathcal R=0$, distance rescalings make every annulus converge
smoothly to a flat annulus in the Tits cone. Its sphere at infinity
is a smooth constant-curvature-one surface and is round by
Claim~\ref{claim:sphere-at-infinity-round}. Thus the cone is
$\Ar^3$. Bishop--Gromov rigidity then gives
$\Vol B(p,R)=\omega_3R^3$ for all $R$, so $(M,g(t))$ is flat, a
contradiction.
\end{proof}

\begin{claim}[Sphere at infinity is round]\label{claim:sphere-at-infinity-round}
In the zero-asymptotic-curvature case above,
$S_\infty(M,p)$ is isometric to the round $2$-sphere.
\end{claim}
\begin{proof}
\sketch The complement of the cone vertex is orientable because
$M$ is orientable. Hence its constant-positive-curvature link is an
orientable surface, necessarily $S^2$ with its round metric.
\end{proof}

\section{Universal $\kappa$}

\begin{proposition}[Universal $\kappa_0$ for non-round $3$-dimensional $\kappa$-solutions]\label{prop:universal-kappa-non-round}
\uses{cor:compact-asymptotic-soliton-classification, thm:gradient-shrinking-soliton-classification}
\label{kappa0}

There is a universal $\kappa_0>0$ such that every non-round
three-dimensional $\kappa$-solution is a $\kappa_0$-solution.
\end{proposition}
\begin{proof}
\sketch
\uses{cor:compact-asymptotic-soliton-classification,thm:gradient-shrinking-soliton-classification}
A non-round solution has a noncompact asymptotic soliton, hence one
of the three cylindrical models $S^2\times\Ar$,
$\Ar P^2\times\Ar$, or the twisted line bundle over $\Ar P^2$.
Choose $q_k$ with $l_x(q_k,\bar\tau_k)\leq3/2$ and
$\bar\tau_k\to\infty$. On a fixed model ball and time interval the
rescaled flows converge smoothly, while reduced length and volume
are bounded uniformly. Theorem~\ref{THM} therefore supplies a
non-collapsing constant depending only on these three models. Scale
invariance gives $\kappa_0$-non-collapsing up to
$\sqrt{2\bar\tau_k}$ at time zero, and $k\to\infty$ gives all
scales. Shifting any $t\leq0$ to time zero proves the assertion
throughout the flow.
\end{proof}

\section{Asymptotic volume}

For a $\kappa$-solution define
\[
 \mathcal V_\infty(M,g(t))
 =\lim_{r\to\infty}\frac{\Vol_{g(t)}B(p,r)}{r^n}.
\]
Bishop--Gromov monotonicity makes the limit independent of $p$.

\begin{theorem}[Asymptotic volume of a $\kappa$-solution is zero]\label{thm:asymptotic-volume-zero}
\uses{def:ancient-solution, def:kappa-noncollapsed}
\label{asympvol}

Every $\kappa$-solution has identically zero asymptotic volume.
\end{theorem}
\begin{proof}
\sketch
\uses{cor:two-dimensional-gradient-shrinking-soliton,prop:splitting-at-infinity}
Induct on dimension. The two-dimensional case is compact by
Corollary~\ref{cor:two-dimensional-gradient-shrinking-soliton}.
Splitting at infinity gives in dimension $n$ a limit
$(N^{n-1},h(t))\times\Ar$. Product-ball containment makes the
factor a $\kappa/2$-ancient solution. The inductive hypothesis gives
zero asymptotic volume for the factor and hence for the product and
the original flow. The proof uses the splitting results established
here in dimensions two and three; the statement and argument extend
to all dimensions when their higher-dimensional versions are
available.
\end{proof}

\subsection{Volume comparison}

\begin{proposition}[Volume comparison for $\kappa$-non-collapsed balls]\label{prop:volume-comparison-bound}
\uses{thm:asymptotic-volume-zero}
\label{volcomp}

Fix $n$. For every $\nu>0$ there is $A<\infty$ with the following
property. Let $(M_k,g_k(t))$, $-t_k\leq t\leq0$, be Ricci flows with
nonnegative curvature operator, let
$B(p_k,0,r_k)\Subset M_k$, and put $Q_k=R(p_k,0)$. Suppose
\[
 R(q,t)\leq4Q_k
 \quad(q\in B(p_k,0,r_k),\ -t_k\leq t\leq0),
\]
and $t_kQ_k,r_k^2Q_k\to\infty$. Then, for all sufficiently large
$k$,
\[
 \Vol B(p_k,0,AQ_k^{-1/2})
 <\nu(AQ_k^{-1/2})^n.
\]
\end{proposition}
\begin{proof}
\sketch
\uses{thm:asymptotic-volume-zero}
Otherwise diagonalize so that the reverse volume inequality holds
for every fixed $A$. Rescaling by $Q_k$ gives a complete nonflat
ancient limit with $R\leq4$, $R(p_\infty,0)=1$, and asymptotic
volume at least $\nu$. Claim~\ref{claim:asymptotic-volume-monotone}
makes the limit a $\nu$-solution, contradicting
Theorem~\ref{thm:asymptotic-volume-zero}.
\end{proof}

\begin{claim}[Monotonicity of the asymptotic volume]\label{claim:asymptotic-volume-monotone}
Let $(M,g(t))$ be an ancient flow of complete manifolds with bounded
nonnegative curvature operator, and let $\mathcal V(t)$ denote its
asymptotic volume.
\begin{enumerate}
\item[(1)] $\mathcal V(t)$ is non-increasing in $t$.
\item[(2)] If $\mathcal V(t)=V>0$, then
$\Vol_{g(t)}B(x,r)\geq Vr^n$ for all $x,r$.
\end{enumerate}
\end{claim}
\begin{proof}
\sketch For $a<b\leq0$, bounded curvature and the distance
distortion estimate give
$B_{g(b)}(p,r)\subset B_{g(a)}(p,r+A(b-a))$. Volume of a fixed set
is non-increasing under nonnegative Ricci flow, so division by
$r^n$ and $r\to\infty$ gives
$\mathcal V(b)\leq\mathcal V(a)$. The second assertion is
Bishop--Gromov comparison.
\end{proof}

\begin{corollary}[Volume bounded below implies curvature bounded above]\label{cor:volume-bound-implies-curvature-bound}
\uses{prop:volume-comparison-bound}
\label{bdvolbdcurv}

For every $\nu>0$ there is $C=C(n,\nu)$ such that, if an ancient
complete flow has bounded nonnegative curvature operator and
\[
 \Vol B(p,0,r)\geq\nu r^n,
\]
then $r^2R(q,0)\leq C$ for every $q\in B(p,0,r)$.
\end{corollary}
\begin{proof}
\sketch
\uses{prop:volume-comparison-bound,lem:point-picking-general,thm:asymptotic-volume-zero}
For a counterexample, point-picking supplies $q_k'$, scales $s_k$,
and $Q_k'=R(q_k',0)$ with $Q_k's_k^2\to\infty$ and
$R<4Q_k'$ on $B(q_k',s_k)$. Bishop--Gromov transfers the original
volume lower bound to every smaller centered ball. After rescaling
by $Q_k'$, a subsequence converges to a nonflat ancient solution
with positive asymptotic volume, contradicting
Theorem~\ref{thm:asymptotic-volume-zero}.
\end{proof}

\begin{corollary}[Global curvature bound from a volume lower bound]\label{cor:global-curvature-bound-from-volume}
\uses{cor:volume-bound-implies-curvature-bound}
\label{global}

For fixed $n$ and $\nu>0$ there is a finite function $K(A)$ on
$(0,\infty)$ such that, under the hypotheses of the preceding
corollary,
\[
 (r+d_0(p,q))^2R(q,0)
 \leq K(d_0(p,q)/r)
\]
for every $q\in M$.
\end{corollary}
\begin{proof}
\sketch
\uses{cor:volume-bound-implies-curvature-bound}
With $d=d_0(p,q)$,
\[
 \Vol B(q,0,r+d)\geq\nu r^n
 =\frac{\nu}{(1+d/r)^n}(r+d)^n.
\]
Apply the preceding corollary at $q$ with radius $r+d$ and
non-collapsing constant $\nu/(1+d/r)^n$.
\end{proof}

\section{Compactness of the space of $3$-dimensional
$\kappa$-solutions}

\begin{theorem}[Compactness of based $3$-dimensional $\kappa$-solutions]\label{thm:kappa-solution-compactness}
\uses{def:ancient-solution, def:kappa-noncollapsed}
\label{compactkappa}

Every sequence of based three-dimensional $\kappa$-solutions
$(M_k,g_k(t),(p_k,0))$ with $R(p_k,0)=1$ has a smoothly convergent
subsequence whose limit is a based $\kappa$-solution.
\end{theorem}
\begin{proof}
\sketch
\uses{lem:uniform-curvature-bound-kappa-solution,cor:three-dimensional-asymptotic-soliton-bounded-curvature}
Lemma~\ref{lem:uniform-curvature-bound-kappa-solution} and
$\partial_tR\geq0$ give uniform curvature bounds on every fixed
based ball for all earlier times. Non-collapsing and flow compactness
give an ancient smooth limit. It is complete, non-flat because
$R(p_\infty,0)=1$, nonnegatively curved, and $\kappa$-non-collapsed.
Corollary~\ref{cor:three-dimensional-asymptotic-soliton-bounded-curvature}
supplies bounded curvature on each limiting slice, so the limit is
a $\kappa$-solution.
\end{proof}

\begin{lemma}[Uniform curvature bound near the base point]\label{lem:uniform-curvature-bound-kappa-solution}
\uses{thm:kappa-solution-compactness}
\label{proposition}

For each finite $r$ there is $C(r)<\infty$ such that every based
three-dimensional $\kappa$-solution with $R(p,0)=1$ satisfies
$R(q,0)\leq C(r)$ on $B(p,0,r)$.
\end{lemma}
\begin{proof}
  \uses{prop:asymptotic-scalar-curvature-infinite, cor:volume-bound-implies-curvature-bound, cor:three-dimensional-asymptotic-soliton-bounded-curvature}
\sketch

Choose the closest $q$ to $p$ with
$d_0(p,q)^2R(q,0)=1$, and write $d=d_0(p,q)$,
$Q=R(q,0)$. Claims~\ref{claim:curvature-ratio-bound-near-basepoint}
and~\ref{claim:curvature-upper-bound-universal} give a universal
upper bound for $Q$ and for $R/Q$ on $B(q,2d)$.
Claim~\ref{claim:basepoint-neighborhood-curvature-bound} therefore
provides a universal ball about $p$ with uniform curvature bound.
Non-collapsing gives a fixed volume lower bound there, and
Corollary~\ref{cor:global-curvature-bound-from-volume} propagates
the curvature estimate to every fixed ball.
\end{proof}

\begin{claim}[Universal curvature ratio bound near $q$]\label{claim:curvature-ratio-bound-near-basepoint}
\uses{cor:volume-bound-implies-curvature-bound, cor:three-dimensional-asymptotic-soliton-bounded-curvature}
\label{2Rbound}

There is a universal $C$ such that
$R(q',0)/R(q,0)\leq C$ for all $q'\in B(q,0,2d)$.
\end{claim}
\begin{proof}
\sketch If no such bound existed, Corollary~\ref{cor:volume-bound-implies-curvature-bound}
would force
\begin{equation}\label{V3}
 \Vol B(q_k,0,2d_k)<\nu(2d_k)^3
\end{equation}
for every fixed $\nu>0$ and large $k$; after diagonalization,
\begin{equation}\label{lowerbd}
 \Vol B(q_k,0,2d_k)<\nu(2d_k)^3
\end{equation}
holds along the sequence. Choose $r_k<2d_k$ with
\begin{equation}\label{V/2}
 \Vol B(q_k,0,r_k)=\frac{\omega_3}{2}r_k^3.
\end{equation}
Then $r_k/d_k\to0$. Rescaling by $r_k^{-2}$ and using global
curvature control gives a complete $\kappa$-non-collapsed ancient
limit with
\begin{equation}\label{V3/2}
 \Vol B(q_\infty,0,1)=\frac{\omega_3}{2}
\end{equation}
and $R(q_\infty,0)=0$. The strong maximum principle makes the limit
flat. It is not $\Ar^3$ by the volume identity, while every
nontrivial complete flat quotient has zero asymptotic volume and is
not non-collapsed on all scales, a contradiction.
\end{proof}

\begin{claim}[Distance bound under rescaling]\label{claim:distance-bound-under-rescaling}
For each $c>0$ there is a universal $\widetilde C(c)$ such that
\[
 d_{g(-cQ^{-1})}(p,q)\leq\widetilde C(c)Q^{-1/2}.
\]
\end{claim}
\begin{proof}
\sketch The preceding ratio bound and $\partial_tR\geq0$ give
$|\Ric|\leq CQ$ along a $g(0)$-minimizing geodesic from $p$ to
$q$ for all earlier times. Its length $L(t)$ satisfies
$L'(t)\geq-CQL(t)$. Since $d=Q^{-1/2}$, integration to
$-cQ^{-1}$ gives the result with $\widetilde C=e^{cC}$.
\end{proof}

\begin{claim}[Universal upper bound for $Q$]\label{claim:curvature-upper-bound-universal}
\uses{claim:distance-bound-under-rescaling}
\label{zkbound}

The quantity $Q=R(q,0)$ is bounded above by a universal constant.
\end{claim}
\begin{proof}
\sketch
\uses{claim:distance-bound-under-rescaling}
The integrated Harnack inequality and the distance claim give
$R(q,-cQ^{-1})\leq\exp(\widetilde C(c)^2/(2c))$ because
$R(p,0)=1$. After rescaling by $Q$, Shi estimates on $B(q,2)$ imply
$\partial_tR(q,t)\leq C''Q^2$ for $-Q^{-1}<t\leq0$. Therefore
\[
 Q\leq cC''Q+\exp(\widetilde C(c)^2/(2c)).
\]
Taking $c=(2C'')^{-1}$ gives a universal upper bound.
\end{proof}

\begin{claim}[Universal neighborhood curvature bound]\label{claim:basepoint-neighborhood-curvature-bound}
\uses{claim:curvature-ratio-bound-near-basepoint, claim:curvature-upper-bound-universal}
\label{bdnearxk}

There are universal $\delta>0$ and $C_1<\infty$ such that
$d(p,q)\geq\delta$ and $R(q',t)\leq C_1$ for
$q'\in B(p,0,d)$ and $t\leq0$.
\end{claim}
\begin{proof}
\sketch Since $d^2Q=1$ and $Q$ is universally bounded,
$d\geq\delta$. The inclusion $B(p,d)\subset B(q,2d)$ and the
curvature-ratio bound give the time-zero estimate; $\partial_tR\geq0$
extends it backward.
\end{proof}

\begin{corollary}[Curvature bound on a fixed universal ball]\label{cor:curvature-bound-fixed-ball}
\uses{claim:basepoint-neighborhood-curvature-bound}
\label{rho}

After reducing the universal $\delta>0$ if necessary,
$R(q',t)\leq\delta^{-2}$ on
$B(p,0,\delta)\times(-\infty,0]$.
\end{corollary}

\begin{corollary}[Universal scale-invariant gradient and time-derivative bounds]\label{cor:curvature-derivative-bounds-universal}
\uses{thm:kappa-solution-compactness}
\label{derivcor}

For each $\kappa>0$ there is $C<\infty$ such that every
three-dimensional $\kappa$-solution satisfies
\begin{equation}\label{gradest}
 \sup_{(x,t)}\frac{|\nabla R(x,t)|}{R(x,t)^{3/2}}<C,
\end{equation}
\begin{equation}\label{timeest}
 \sup_{(x,t)}\frac{|\partial_tR(x,t)|}{R(x,t)^2}<C.
\end{equation}
The constant may be chosen independently of $\kappa$.
\end{corollary}
\begin{proof}
\sketch
\uses{thm:kappa-solution-compactness}
Both ratios are scale invariant and uniformly bounded on the compact
space of normalized based $\kappa$-solutions. The universal
non-collapsing proposition treats every non-round solution; the
round family satisfies the same bounds directly. Using
Equation~\ref{Revol}, the second bound is equivalently a bound for
$|\triangle R+2|\Ric|^2|/R^2$.
\end{proof}

\section{Qualitative description of $\kappa$-solutions}

\subsection{Strong canonical neighborhoods}

\begin{definition}[$(C,\epsilon)$-cap]\label{def:c-epsilon-cap}
\uses{def:epsilon-neck}
\label{epscap}

Fix $0<\epsilon<1/2$ and $C<\infty$. A
\emph{$(C,\epsilon)$-cap} in a Riemannian $3$-manifold is an open
submanifold $\mathcal C$, together with an open $N\subset\mathcal C$,
such that:
\begin{enumerate}
\item[(1)] $\mathcal C$ is diffeomorphic to an open $3$-ball or a
punctured $\Ar P^3$;
\item[(2)] $N$ is an $\epsilon$-neck and
$\mathcal C\setminus N$ is compact;
\item[(3)] $\overline Y=\mathcal C\setminus N$ is a compact
submanifold with boundary, its interior $Y$ is the \emph{core}, and
$\partial\overline Y$ is a central $2$-sphere of an
$\epsilon$-neck in $\mathcal C$;
\item[(4)] $R>0$ on $\mathcal C$ and
\[
 \diam(\mathcal C)<C\left(\sup_{\mathcal C}R\right)^{-1/2};
\]
\item[(5)] $\sup_{x,y\in\mathcal C}R(y)/R(x)<C$;
\item[(6)]
$\Vol\mathcal C<C(\sup_{\mathcal C}R)^{-3/2}$;
\item[(7)] for $y\in Y$, let $r_y$ satisfy
$\sup_{B(y,r_y)}R=r_y^{-2}$. Then
$B(y,r_y)\Subset\mathcal C$ and
\[
 C^{-1}<\inf_{y\in Y}\frac{\Vol B(y,r_y)}{r_y^3};
\]
\item[(8)]
\[
 \sup_{\mathcal C}\frac{|\nabla R|}{R^{3/2}}<C,
 \qquad
 \sup_{\mathcal C}\frac{|\triangle R+2|\Ric|^2|}{R^2}<C.
\]
\end{enumerate}
\end{definition}

\begin{remark}[The ball $B(y,r_y)$ has compact closure in the cap]\label{rem:cap-core-ball-remark}
\uses{def:c-epsilon-cap}

If $B(y,r_y)$ meets the complement of the core, curvature control on
the neck bounds $r_y$ by a fixed multiple of the neck curvature
scale. Since $\epsilon<1/2$ and $y$ lies in the core, this implies
$B(y,r_y)\Subset\mathcal C$. The neck orientation is chosen so that
the closure of its negative end meets the core.
\end{remark}

\begin{claim}[Condition (8) is a time-derivative bound]\label{claim:cap-condition-time-derivative-equivalence}
\uses{def:c-epsilon-cap, def:ricci-flow}
\label{partialtR}

When $\mathcal C$ lies in a time-slice of a Ricci flow, the second
inequality in Condition~(8) is equivalent to
\[
 \sup_{\mathcal C}\frac{|\partial_tR|}{R^2}<C.
\]
\end{claim}
\begin{proof}
\sketch Use $\partial_tR=\triangle R+2|\Ric|^2$.
\end{proof}

\begin{definition}[$C$-component]\label{def:c-component}
\label{Ccomponent}
A compact connected Riemannian $3$-manifold $(M,g)$ is a
\emph{$C$-component} if:
\begin{enumerate}
\item[(1)] $M$ is diffeomorphic to $S^3$ or $\Ar P^3$;
\item[(2)] $M$ has positive sectional curvature;
\item[(3)]
\[
 C^{-1}<\frac{\inf_PK(P)}{\sup_MR};
\]
\item[(4)]
\[
 C^{-1}\sup_MR^{-1/2}<\diam M<C\inf_MR^{-1/2}.
\]
\end{enumerate}
\end{definition}

\begin{definition}[$\epsilon$-round component]\label{def:epsilon-round-component}
\label{defnepsround}
A compact connected $3$-manifold $(M,g)$ is
\emph{within $\epsilon$ of round in the $C^{[1/\epsilon]}$ topology}
if there are $R>0$, a compact constant-curvature-one manifold
$(Z,g_0)$, and a diffeomorphism $\varphi:Z\to M$ such that
$\varphi^*(Rg)$ is within $\epsilon$ of $g_0$ in that topology.
\end{definition}

\begin{definition}[$(C,\epsilon)$-canonical neighborhood]\label{def:canonical-neighborhood}
\uses{def:epsilon-neck, def:c-epsilon-cap, def:c-component, def:epsilon-round-component}

An open neighborhood $U$ of $x$ in a Riemannian manifold is a
\emph{$(C,\epsilon)$-canonical neighborhood} if it is one of:
\begin{enumerate}
\item[(1)] an $\epsilon$-neck centered at $x$;
\item[(2)] a $(C,\epsilon)$-cap whose core contains $x$;
\item[(3)] a $C$-component;
\item[(4)] an $\epsilon$-round component.
\end{enumerate}
\end{definition}

\begin{definition}[Evolving and strong $\epsilon$-neck; strong canonical neighborhood]\label{def:strong-epsilon-neck-canonical-neighborhood}
\uses{def:generalized-ricci-flow, def:epsilon-neck, def:c-epsilon-cap, def:c-component, def:epsilon-round-component, def:epsilon-close-topology}
\label{defncanonnbhd}

Let $(\mathcal M,G)$ be a generalized Ricci flow, let
$x\in\mathcal M$ with $\mathbf t(x)=t$, and fix $C<\infty$ and
$\epsilon>0$. An \emph{evolving $\epsilon$-neck of normalized time
length $t'>0$ centered at $x$} is an embedding
\[
 \psi:S^2\times(-\epsilon^{-1},\epsilon^{-1})\longrightarrow
 N\subset M_t,qquad x\in\psi(S^2\times\{0\}),
\]
such that:
\begin{enumerate}
\item[(1)] there is a time- and vector-field-compatible embedding
$N\times(t-R(x)^{-1}t',t]\to\mathcal M$;
\item[(2)] after pulling back by $\psi$ and scaling by $R(x)$, the
resulting metric family is within $\epsilon$ in the
$C^{[1/\epsilon]}$ topology of
$(h(s),ds^2)$ for $-t'<s\leq0$, where $h(s)$ is round with scalar
curvature $(1-s)^{-1}$.
\end{enumerate}
A \emph{strong $\epsilon$-neck} has normalized time length $1$.

A \emph{strong $(C,\epsilon)$-canonical neighborhood} of $x$ is an
open $U\subset M_t$ which is one of:
\begin{enumerate}
\item[(1)] a strong $\epsilon$-neck centered at $x$;
\item[(2)] a $(C,\epsilon)$-cap whose core contains $x$;
\item[(3)] a $C$-component;
\item[(4)] an $\epsilon$-round component.
\end{enumerate}
All these notions are invariant under simultaneous parabolic
scaling.
\end{definition}

\begin{proposition}[Continuity of canonical neighborhoods]\label{prop:canonical-neighborhood-stability}
\uses{def:strong-epsilon-neck-canonical-neighborhood, def:canonical-neighborhood}
\label{canonvary}

Fix $\epsilon<1/4$ and $C<\infty$.
\begin{enumerate}
\item[(1)] If $U\subset M_t$ is a $(C,\epsilon)$-canonical
neighborhood of $x$, then every horizontal metric sufficiently
close to $G|_U$ in $C^{[1/\epsilon]}$ makes $U$ a
$(C,\epsilon)$-canonical neighborhood of every sufficiently nearby
$x'$.
\item[(2)] If an evolving $\epsilon$-neck centered at $(x,t)$ has
normalized time length $a>1$, every sufficiently close Ricci flow on
the same spacetime cylinder contains a strong $\epsilon$-neck
centered at $(x,t)$.
\item[(3)] If $C'>C$, $\epsilon'>\epsilon$, and $R(x)\leq2$, there
is $\delta>0$ such that every metric $\delta$-close to a
$(C,\epsilon)$-canonical neighborhood contains a
$(C',\epsilon')$-canonical neighborhood of $x$.
\item[(4)] If $g(s)$, $-1<s\leq0$, is a strong $\epsilon$-neck
centered at $(x,0)$ with $R_g(x,0)=1$, then every sufficiently close
normalized family is a strong $\epsilon'$-neck.
\end{enumerate}
\end{proposition}
\begin{proof}
\sketch Diameter, volume, curvature, curvature ratios, and the
derivative quantities in the definitions vary continuously in the
stated topology. The same neck chart persists under small metric
perturbations and, after precomposition with a small axial
diffeomorphism, under small changes of center. These observations
prove (1) for necks and components; for caps they apply to both neck
charts and Conditions (4)--(8), including continuity of $r_y$.
An interval of normalized length strictly greater than one absorbs
small curvature-scale changes, proving (2). For (3), strict
inequalities absorb the changes $C<C'$ and
$\epsilon<\epsilon'$; truncate and recenter the end neck while
keeping the same cap core. Statement (4) is the corresponding
spacetime continuity assertion.
\end{proof}

\begin{corollary}[Centers of strong necks form an open set]\label{cor:strong-neck-centers-open}
\uses{prop:canonical-neighborhood-stability}
\label{strepsopen}

In an ancient solution, the set of spacetime points that are centers
of strong $\epsilon$-necks is open.
\end{corollary}
\begin{proof}
\sketch
\uses{prop:canonical-neighborhood-stability}
A strong neck in an ancient flow extends slightly beyond normalized
backward time one. Nearby times still have an evolving neck of
length greater than one, and a small axial diffeomorphism moves its
center to any nearby spatial point. Apply Part~(2) of the
proposition.
\end{proof}

\begin{definition}[$\epsilon$-tube and capped $\epsilon$-tube]\label{def:epsilon-tube}
\uses{def:epsilon-neck, def:c-epsilon-cap}
\label{epstube}

An \emph{$\epsilon$-tube} is a submanifold diffeomorphic to
$S^2\times I$, for a nondegenerate interval $I$, such that every
boundary component is the central sphere of an $\epsilon$-neck and
the tube is the union of $\epsilon$-necks and the closed half-necks
at its boundary. Every central sphere is isotopic in the tube to an
$S^2$ factor. An \emph{open $\epsilon$-tube} has no boundary and is
a union of such necks.

A \emph{$C$-capped $\epsilon$-tube} is the connected union of a
$(C,\epsilon)$-cap and an open $\epsilon$-tube whose intersection is
$S^2\times(0,1)$ and contains an end of each. A
\emph{doubly $C$-capped $\epsilon$-tube} is the closed connected
union of an open tube and two $(C,\epsilon)$-caps with disjoint core
closures attached at opposite ends.

An \emph{$\epsilon$-fibration} is a closed connected
$S^2$-bundle over $S^1$ that is a union of $\epsilon$-necks whose
central spheres are isotopic to the fibers.
\end{definition}

\begin{definition}[Strong $\epsilon$-tube]\label{def:strong-epsilon-tube}
\uses{def:epsilon-tube, def:strong-epsilon-neck-canonical-neighborhood}

A \emph{strong $\epsilon$-tube} in a generalized Ricci flow is an
$\epsilon$-tube every point of which is the center of a strong
$\epsilon$-neck.
\end{definition}

\subsection{Canonical neighborhoods for $\kappa$-solutions}

\begin{proposition}[Trichotomy for $3$-dimensional $\kappa$-solutions]\label{prop:kappa-solution-trichotomy}
\uses{def:ancient-solution, def:kappa-noncollapsed}
\label{infprod2}

Every three-dimensional $\kappa$-solution is exactly one of:
\begin{enumerate}
\item[(1)] positively curved at every time;
\item[(2)] the product of shrinking round $S^2$'s with a line;
\item[(3)] a flow on a line bundle over $\Ar P^2$ having a finite
cover of type (2).
\end{enumerate}
\end{proposition}
\begin{proof}
\sketch If positive curvature fails, the strong maximum principle
gives a one- or two-sheeted cover splitting into a round surface
flow and a flat circle or line. A fixed circle factor violates
non-collapsing as $t\to-\infty$. The remaining factor is a line; a
further double cover replaces $\Ar P^2$ by $S^2$, yielding (2) or
(3).
\end{proof}

\begin{lemma}[Neck structure outside a bounded ball]\label{lem:noncompact-kappa-solution-neck-structure}
\uses{prop:kappa-solution-trichotomy}

Let $(M,g(t))$ be a noncompact positively curved
three-dimensional $\kappa$-solution and $p\in M$. There are
$D',D_1'<\infty$, possibly depending on the flow and $p$, such that:
every point outside
$B(p,0,D'R(p,0)^{-1/2})$ is the center of an evolving
$\epsilon$-neck of normalized time length $2$; and for
$x$ in this ball and every $2$-plane $P_x$,
\[
 (D_1')^{-1}<\frac{K(P_x)}{R(p,0)}<D_1'.
\]
\end{lemma}
\begin{proof}
\sketch
\uses{thm:kappa-solution-compactness,cor:two-dimensional-gradient-shrinking-soliton,lem:uniform-curvature-bound-kappa-solution}
Normalize $R(p,0)=1$ and suppose non-neck points $p_k$ escape to
infinity. If $d(p,p_k)^2R(p_k,0)\to\infty$, compactness and splitting
at infinity yield the round cylinder, so $p_k$ are eventually
evolving neck centers. If that product remains bounded, rescaling by
$R(p_k,0)\to0$ places curvature $R(p,0)/R(p_k,0)\to\infty$ at
bounded based distance, contradicting the uniform local curvature
bound. The curvature comparison on the remaining compact ball
follows from strict positive curvature.
\end{proof}

\begin{proposition}[Uniform neck-structure constants]\label{prop:uniform-neck-structure-constants}
\uses{lem:noncompact-kappa-solution-neck-structure}
\label{uniformD}

For every sufficiently small $\epsilon>0$ there are universal
$D(\epsilon),D_1(\epsilon)<\infty$ such that, for any noncompact
positively curved three-dimensional $\kappa$-solution and any soul
$p$ of $(M,g(0))$:
\begin{enumerate}
\item[(1)] every point outside
$B(p,0,DR(p,0)^{-1/2})$ is the center of a strong
$\epsilon$-neck. Inside this ball,
\[
 D_1^{-1}<\frac{K(P_x)}{R(p,0)}<D_1,
\]
and
\[
 D_1^{-3/2}R(p,0)^{-3/2}
 <\Vol B(p,0,DR(p,0)^{-1/2})
 <D_1^{3/2}R(p,0)^{-3/2}.
\]
\item[(2)] The soul misses the middle two-thirds of every
$\epsilon$-neck. Every central neck sphere is isotopic in
$M\setminus\{p\}$ to every distance sphere $d(p,\cdot)^{-1}(a)$,
$a>0$. The region between two disjoint central neck spheres is
diffeomorphic to $S^2\times[0,1]$.
\end{enumerate}
\end{proposition}
\begin{proof}
\sketch
\uses{thm:kappa-solution-compactness,lem:noncompact-kappa-solution-neck-structure,prop:canonical-neighborhood-stability,prop:kappa-solution-trichotomy,cor:strong-neck-centers-open}
If no universal $D$ existed, normalize souls $p_k$ by
$R(p_k,0)=1$ and take escaping non-neck points. Compactness gives a
noncompact limit whose base soul is not a neck center. The trichotomy
therefore makes the limit positively curved, and the preceding lemma
gives necks outside a fixed limiting ball. Stability transfers these
necks back to larger and larger annuli. Choosing the first
non-neck point $q_k$ beyond such an annulus puts it in the closure of
neck centers and hence in a $2\epsilon$-neck.

Let $Q_k=R(q_k,0)$. Claim~\ref{claim:kq-scale-separation} shows that
the rescaled distance to $p_k$ diverges. Compactness and splitting
then make the based $Q_k$-rescaled limit a round cylinder, so $q_k$
was a strong $\epsilon$-neck center, a contradiction. Compactness at
the soul gives the universal curvature and volume constants $D_1$.
The disjointness and isotopy assertions are Lemma~\ref{topiso} and
Corollary~\ref{neckseparate}.
\end{proof}

\begin{remark}[Basepoint may be replaced]\label{rem:uniform-neck-structure-basepoint-remark}
\uses{prop:uniform-neck-structure-constants}

In Part~(1), the soul may be replaced by any point that is not the
center of a strong $\epsilon$-neck.
\end{remark}

\begin{claim}[Curvature scale separates $p_k$ and $q_k$]\label{claim:kq-scale-separation}
In the contradiction sequence above,
\[
 d_0(p_k,q_k)^2Q_k\longrightarrow\infty.
\]
\end{claim}
\begin{proof}
\sketch
\uses{lem:uniform-curvature-bound-kappa-solution}
If the products were bounded, then $d_0(p_k,q_k)\to\infty$ would
force $Q_k\to0$. In the $Q_k$-rescaled flows, $p_k$ stays at bounded
distance from $q_k$ but has scalar curvature $Q_k^{-1}\to\infty$,
contradicting the uniform curvature bound near a normalized base
point.
\end{proof}

\begin{corollary}[Non-compact $\kappa$-solutions are (capped) strong tubes]\label{cor:noncompact-kappa-solution-tube-structure}
\uses{prop:uniform-neck-structure-constants, cor:curvature-derivative-bounds-universal, def:strong-epsilon-tube}
\label{noncompkappa}

There is $\bar\epsilon_2>0$ such that, for
$0<\epsilon\leq\bar\epsilon_2$, some $C_0(\epsilon)$ has the
following property: every noncompact three-dimensional
$\kappa$-solution containing no embedded $\Ar P^2$ with trivial
normal bundle has time-zero slice either a strong $\epsilon$-tube
or a $C_0$-capped strong $\epsilon$-tube.
\end{corollary}
\begin{proof}
\sketch
\uses{prop:uniform-neck-structure-constants,cor:curvature-derivative-bounds-universal}
Take $C_0$ larger than the uniform neck, core-curvature, and
derivative constants. A positively curved solution has the required
cap and neck decomposition. The product cylinder is entirely a
strong tube. For its nontrivial orientable quotient, the invariant
central sphere has a controlled cap neighborhood, and all points at
distance at least $3\epsilon^{-1}$ from its quotient are strong neck
centers.
\end{proof}

\begin{theorem}[Structure of compact $\kappa$-solutions]\label{thm:compact-kappa-solution-structure}
\uses{cor:noncompact-kappa-solution-tube-structure, def:strong-epsilon-tube}
\label{thmep3}

There is $\bar\epsilon_3>0$ such that, for every
$0<\epsilon\leq\bar\epsilon_3$, some $C_1(\epsilon)<\infty$ makes
every compact three-dimensional $\kappa$-solution satisfy one of:
\begin{enumerate}
\item[(1)] $M$ has constant positive sectional curvature;
\item[(2)]
$\diam(M,g(0))<C_1(\max_MR(\cdot,0))^{-1/2}$ and $M$ is
diffeomorphic to $S^3$ or $\Ar P^3$;
\item[(3)] $(M,g(0))$ is a doubly $C_1$-capped strong
$\epsilon$-tube.
\end{enumerate}
\end{theorem}
\begin{proof}
  \uses{prop:universal-kappa-non-round, cor:noncompact-kappa-solution-tube-structure, cor:compact-asymptotic-soliton-classification, thm:gradient-shrinking-soliton-classification, thm:kappa-solution-compactness, prop:canonical-neighborhood-stability}
\sketch

A compact $\kappa$-solution has positive curvature: otherwise its
universal cover splits as $S^2\times\Ar$, forcing noncompactness or
a collapsed $S^2\times S^1$ quotient. Its fundamental group is
finite, so it contains no embedded $\Ar P^2$ with trivial normal
bundle. For a non-round solution, universal non-collapsing and
Claim~\ref{claim:diameter-canonical-neighborhood-criterion} show
that sufficiently large diameter gives a cap/neck decomposition.
The topology theorem leaves $S^3$, $\Ar P^3$,
$\Ar P^3\#\Ar P^3$, or an $S^2$-bundle over $S^1$; finite
fundamental group excludes the last two \cite{Petersen}, giving (3).
If the diameter is small, the asymptotic soliton is noncompact and
cylindrical; at sufficiently negative time the large-diameter case
applies, showing that $M$ is $S^3$ or $\Ar P^3$.
\end{proof}

\begin{claim}[Large-diameter compact solutions are capped tubes]\label{claim:diameter-canonical-neighborhood-criterion}
\uses{cor:noncompact-kappa-solution-tube-structure}

For each $\epsilon>0$ there is $C_1$ such that a compact non-round
solution with
\[
 \diam(M,g(0))>C_1(\max_MR(\cdot,0))^{-1/2}
\]
has every point either in the core of a
$(C_0(\epsilon),\epsilon)$-cap or at the center of a strong
$\epsilon$-neck.
\end{claim}
\begin{proof}
\sketch
\uses{thm:kappa-solution-compactness,cor:noncompact-kappa-solution-tube-structure,prop:canonical-neighborhood-stability}
Otherwise normalize a violating point to scalar curvature one in a
sequence whose normalized diameters diverge. Compactness yields a
noncompact based limit. Its tube decomposition puts the limiting
point in a cap core or strong neck, and canonical-neighborhood
stability transfers that structure to the violating points.
\end{proof}

\begin{proposition}[Volume and curvature bounds for small-diameter compact solutions]\label{prop:compact-kappa-solution-volume-curvature-bounds}
\uses{thm:compact-kappa-solution-structure, prop:universal-kappa-non-round, thm:kappa-solution-compactness, cor:noncompact-kappa-solution-tube-structure}

Let $0<\epsilon\leq\min(\bar\epsilon_2,\bar\epsilon_3)$ and let
$C_1(\epsilon)$ be as above. There is $C_2(\epsilon)<\infty$ such
that, for every compact non-round $\kappa$-solution satisfying
\[
 \diam(M,g(0))<C_1(\max_MR(\cdot,0))^{-1/2},
\]
every $x,y\in M$ and $2$-plane $P_y\subset T_yM$ satisfy
\[
 C_2^{-1}R(x,0)^{-3/2}<\Vol(M,g(0))
 <C_2R(x,0)^{-3/2},
\]
\[
 C_2^{-1}<\frac{K(P_y)}{R(x,0)}<C_2.
\]
\end{proposition}
\begin{proof}
\sketch
\uses{prop:universal-kappa-non-round,thm:kappa-solution-compactness}
Universal non-collapsing reduces the normalized family to the compact
space in Theorem~\ref{thm:kappa-solution-compactness}; the displayed
scale-invariant quantities have positive finite extrema there.
\end{proof}

\begin{remark}[No universal volume lower bound for round solutions]\label{rem:round-kappa-solution-volume-remark}
For a round $\kappa$-solution, scalar curvature is spatially
constant and volume is bounded above by a universal multiple of
$R^{-3/2}$. A lower bound must contain the factor
$|\pi_1(M)|^{-1}$; it has the form
$C_2|\pi_1(M)|^{-1}R^{-3/2}$.
\end{remark}

\begin{theorem}[Summary structure theorem for $\kappa$-solutions]\label{thm:kappa-solution-structure-summary}
\uses{thm:compact-kappa-solution-structure, cor:noncompact-kappa-solution-tube-structure, prop:compact-kappa-solution-volume-curvature-bounds, cor:curvature-derivative-bounds-universal}
\label{kappasummary}

There is $\bar\epsilon>0$ such that, for
$0<\epsilon<\bar\epsilon$, some $C=C(\epsilon)$ makes every
three-dimensional $\kappa$-solution satisfy exactly one of:
\begin{enumerate}
\item[(1)] it is round for all $t\leq0$, and $M=S^3/\Gamma$ for a
finite freely acting $\Gamma<SO(4)$;
\item[(2)] $(M,g(0))$ is compact and positively curved,
$M\cong S^3$ or $\Ar P^3$, and for all $x,y\in M$ and
$P_y\subset T_yM$,
\begin{eqnarray*}
 C^{-1/2}R(x,0)^{-1/2}&<&\diam(M,g(0))<CR(x,0)^{-1/2},\\
 C^{-1}R(x,0)^{-3/2}&<&\Vol(M,g(0))<CR(x,0)^{-3/2},\\
 C^{-1}R(x,0)&<&K(P_y)<CR(x,0);
\end{eqnarray*}
\item[(3)] $(M,g(0))$ is positively curved and a doubly
$C$-capped strong $\epsilon$-tube, with
$M\cong S^3$ or $\Ar P^3$;
\item[(4)] $(M,g(0))$ is positively curved and a $C$-capped strong
$\epsilon$-tube, with $M\cong\Ar^3$;
\item[(5)] $(M,g(0))$ is the quotient of round
$S^2\times\Ar$ by a free orientation-preserving involution; it is a
$C$-capped strong $\epsilon$-tube and a punctured $\Ar P^3$;
\item[(6)] $(M,g(0))$ is round $S^2\times\Ar$ and is a strong
$\epsilon$-tube;
\item[(7)] $(M,g(0))$ is $\Ar P^2\times\Ar$, with the
$\Ar P^2$ factor of constant Gaussian curvature.
\end{enumerate}
Except in (1), (2), and (7), every point is in a cap core or is a
strong neck center. In every case,
\begin{equation}\label{nablaR}
 \sup_{p\in M,t\leq0}\frac{|\nabla R(p,t)|}{R(p,t)^{3/2}}<C,
\end{equation}
\begin{equation}\label{dRdt}
 \sup_{p\in M,t\leq0}\frac{|\partial_tR(p,t)|}{R(p,t)^2}<C.
\end{equation}
\end{theorem}

\begin{corollary}[Existence of canonical neighborhoods]\label{cor:canonical-neighborhood-existence}
\uses{thm:kappa-solution-structure-summary, def:canonical-neighborhood}
\label{kappacannbhd}

For every $0<\epsilon\leq\bar\epsilon'$ there is
$C(\epsilon)<\infty$ such that every point of a $\kappa$-solution
has a strong $(C,\epsilon)$-canonical neighborhood, unless the
solution is $\Ar P^2\times\Ar$.
\end{corollary}

\begin{corollary}[Canonical neighborhoods pass to limits of generalized flows]\label{cor:canonical-neighborhood-limit-passage}
\uses{cor:canonical-neighborhood-existence, def:strong-epsilon-neck-canonical-neighborhood, def:generalized-ricci-flow}
\label{limitcannbhd}

Fix $0<\epsilon\leq\bar\epsilon'$ and the preceding $C(\epsilon)$.
Let $(\mathcal M_n,G_n,x_n)$ be based generalized Ricci flows with
$\mathbf t(x_n)=0$, no time-slice containing an embedded
$\Ar P^2$ with trivial normal bundle, and a smooth based limit
$(M_\infty,g_\infty(t),(x_\infty,0))$, $-\infty<t\leq0$, which is
a $\kappa$-solution. Then $x_n$ has a strong
$(C,\epsilon)$-canonical neighborhood for all sufficiently large
$n$.
\end{corollary}
\begin{proof}
\sketch
\uses{cor:canonical-neighborhood-existence,prop:canonical-neighborhood-stability}
The limit contains no embedded $\Ar P^2$ with trivial normal
bundle, hence $x_\infty$ has a canonical neighborhood. Smooth
convergence transfers round components and $C$-components directly.
A strong neck extends to normalized time length greater than one and
therefore persists by Part~(2) of
Proposition~\ref{prop:canonical-neighborhood-stability}; cap cores
persist by Part~(1).
\end{proof}
